\section{Related Literature}
Our work combines concepts from automated fact-checking, search-based verification, RAG, and LLM-based self-refinement and introduces a novel application of AI in venture capital due diligence. In the following, we will review each area and explain our contributions.
\\
\\
\textbf{Automated fact-checking.} The task of factual claims verification has been studied extensively in NLP. \citet{guo2022survey} provides a comprehensive survey defining a three-stage pipeline: claim detection, evidence retrieval, and veracity prediction. Subsequent work extended verification to structured data \citep{aly2021feverous} and real-world web claims \citep{schlichtkrull2024averitec}. Two recent surveys document the shift toward LLM-based approaches. \citet{dmonte2025claim} examines how RAG and prompting strategies are reshaping claim verification pipelines, while \citet{vykopal2024generative} focuses specifically on LLMs across all stages of fact-checking. Our work follows the general pipeline structure but targets a novel domain -- venture capital due diligence -- where claims are domain-specific, unverifiable through standard knowledge bases, and must be linked to investment risk.
\\
\\
\textbf{Search-based verification.} Several systems have combined LLMs with external tools for factuality assessment. FacTool \citep{chern2023factool} proposes a task-agnostic framework that extracts claims from LLM-generated text, generates search queries, gathers evidence via Google Search and other tools, and verifies each claim. SAFE \citep{wei2024longform} takes a similar approach to evaluating long-form factuality, decomposing responses into atomic facts and verifying each using multi-step Google Search reasoning. OpenFactCheck \citep{wang2024openfactcheck} provides a modular framework that allows claim processors, retrievers, and verifiers to be freely combined and configured. Our pipeline uses search-based verification similar to FacTool and SAFE. However, unlike these systems, which primarily verify LLM-generated text or political claims, \textsc{VentrureTruth} targets structured business documents and connects verification outcomes to the risk assessments.
\\
\\
\textbf{RAG.} RAG has become a dominant paradigm for grounding LLM outputs in external knowledge \citep{gao2023retrieval}. By retrieving relevant documents before generation, RAG systems reduce hallucinations and enable access to the latest information. \citet{adjali2024exploring} applies RAG specifically to the AVeriTeC real-world claim verification shared task, integrating dense retrieval with question generation in an end-to-end setup. In our system, we use Perplexity's Sonar Pro model as a search-augmented retriever that queries the internet, ensuring that verification evidence reflects the most current publicly available information.
\\
\\
\textbf{LLM-based self-refinement.} \citet{madaan2023selfrefine} introduces Self-Refine, demonstrating that LLMs can iteratively improve their outputs through a feedback-refine loop without additional training. \citet{shinn2023reflexion} propose Reflexion, where language agents use verbal self-reflection and dynamic memory to improve task performance over successive attempts. Our pipeline incorporates a similar self-reflection mechanism: after initial claim verification, the system reviews its own outputs for completeness and consistency, iteratively refining verification results.
\\
\\
\textbf{AI in venture capital due diligence.} While commercial tools such as Caena, Tracxn, and AlphaSense have begun integrating AI for deal sourcing, document review, and financial analysis, there is limited academic work specifically addressing automated claim verification in the context of startup evaluation. Existing tools focus primarily on data aggregation and pattern recognition rather than end-to-end factual verification of specific claims made in pitch materials. \textsc{VentureTruth} addresses this gap by providing a transparent, traceable pipeline that extracts individual claims from startup pitch documents, verifies each against web sources, and produces structured risk assessments for investor decision-making.
